\cotefile{83}
\scanpage{113}
pas de l'ordre dans lequel on prend les deux \deleted{\unsure{pai|paires}} couples $(D_1,D_2)$, $(D'_1,D'_2)$. On a un rapport plus intrinsèque, qui ne change pas par permutation $D_1\leftrightarrow D_2$ ou $D'_1\leftrightarrow D'_2$
\[
\rho=\lambda+\lambda^{-1}
\]
qui est défini sur $k$ (sans localiser). On retrouve $\lambda$, au passage à $\lambda^{-1}$ près, comme solution de
\[
\lambda^2-\lambda\rho+1=0\;.
\]

\deleted{\unsure{Ce|Le} $\rho$ est \unsure{part|pour}} On aura ici
\[
\rho\ne 2\qquad\text{sur chaque fibre.}
\]
i.e.
\[
\rho_0=\rho-2\qquad\begin{aligned}&\text{est $\ne0$ sur chaque fibre}\\[-2pt]&\text{i.e. est invers.}\end{aligned}
\]

On voit \deleted{que} aisément que les classes de \unsure{\deleted{pai}res|paires|\deleted{paires}}, \deleted{dans} $(H,H')$ dans $A$, \inserted{($H,H'$ étales sur $k$, $H,H'$ en pos. gén.),} modulo \unsure{conjugaisons|conjugaison} locales correspondent aux diff $\rho_0$ inv. i.e. \unsure{si on|si on a|si} deux telles paires $(H,H')$ avec inv. $\rho_0=\rho_0(H,H')$, $(\overline H,\overline H')$ avec inv. $\overline{\rho_0}$
% archive sheet number: 112
\scanpage{114}
alors $\exists$ \inserted{local\textsuperscript{t}} autom \deleted{\unsure{int.}}\inserted{$\theta$} de $A$, tel que
\[
\theta(H)=\overline H)\qquad\theta(H')=\overline{H'}
\]
ssi
\[
\overline{\rho_0}=\rho_0
\]

D'autre part, si on se donne $\rho_0$ \inserted{(inv)},  il existe \uline{\unsure{loc|alors|local\textsuperscript{t}}} une paire $(H,H')$ qui le réalise.

D'ailleurs, \deleted{tout autom de $A$ \deleted{\unsure{qui}} \inserted{est} \deleted{\illegible{2}} intérieur loc. Zar (et même glob.\textsuperscript{ment}, si $\operatorname{Pic}(k)=0$), et s'il fixe $H$ il \deleted{est loc. \unsure{un|de}} \inserted{provient loc. Zar} d'une section de $H$ ($H$ est son propre centralisateur dans $A$) \deleted{\unsure{ou} est} provient d'un $\theta\in$ \deleted{de la \unsure{façon|forme}} $x\longmapsto x'=\operatorname{Tr}(x)1-x$}

\deleted{La \unsure{conjug.} par un tel autom provient de} Dans $G=\underline{\operatorname{Aut}}(A)\simeq\underline{\operatorname{Aut}}(P)$, le sous-schéma en groupes qui stabilise à la fois $H$ et $H'$ est un sous-\deleted{gr}\inserted{schéma}
% slanted
\margin{du moins si $\rho\ne-2$ \deleted{partout} i.e. $\rho_0\ne-4$ \deleted{partout}, sinon il y a un groupe $\unsure{C|\overline C}$ d'autos qui est $\simeq\mathbb Z/2\mathbb Z\times\mathbb Z/2\mathbb Z$.} en groupes isom. à $(\mathbb Z/2\mathbb Z)_S$, \deleted{\unsure{qui}} dont la section non nulle est un auto- qui, sur $H$ et sur $H'$ \deleted{est donné par}
% archive sheet number: 113
\scanpage{115}
induit sur $H$ et $H'$ l'automorphisme
\[
x\longmapsto x'=\operatorname{Tr}(x)\cdot1-\unsure{x|a}
\]

En tous cas, les isom qui envoient $(H,H')$ sur $(\overline H,\overline H')$ forment un torseur sous le groupe $\mathbb Z/2\mathbb Z$ (si $\rho_0+4$ inv.) lequel torseur peut (je présume) être non trivial (même sachant que le torseur associé de groupe $\underline{\operatorname{Aut}}(A)\simeq\deleted{\unsure{A^*/}}\check V(A)^*/G_m$ est trivial

\deleted{\unsure{On|En}} En effet : [Si $P$ est ce \inserted{$\mathbb Z/2$-}torseur sur $\operatorname{Spec} k$, on aura un hom
\[
\inserted{P\xrightarrow{\sim}\operatorname{isom}(\unsure{\xi|\mathcal E|\varepsilon}_H,\xi_{\overline H})\;(\simeq\xi_H\wedge\xi_{\overline H})}
\]
\[
\deleted{P\times\xi_H\longrightarrow\xi_{\overline H}}
\]
% slanted
\margin{$\xi_H$ désigne le rev. quadratique étale de $S=\operatorname{Spec}k$ défini par $H$} \deleted{qui définit} et de même
\[
P\xrightarrow{\sim}\xi_{H'}\wedge\xi_{\overline H'}\quad]
\]

\deleted{\unsure{Donc}} $P$ n'a de sections que si $\xi_{H\deleted{\unsure{'}}}\simeq\xi_{\overline H}$, \inserted{$\xi_{H'}\simeq\xi_{\overline{H'}}$} ; Prenant par exemple $H'=\overline{H'}$ correspondant à la bisection $\{0,\infty\}$ de $\mathbb P^1$, et $H,\overline H$ étant les bisections définies resp.\textsuperscript{t} par
% slanted
\margin{sur un corps\ldots}
\[
\begin{aligned}
z^2-bz+c&=0 & b^2-4c&\mathrel{\deleted{\unsure{=}}\ne}0\\ z^2-\overline b z+\overline c&=0 & \overline b^2-4\overline c&\ne0
\end{aligned}
\]
% archive sheet number: 114
\scanpage{116}
on a ici
\[
\rho=\frac{b^2}{c}-2,\qquad\overline\rho=\frac{\overline b^2}{\overline c}-2
\]
la condition $\rho=\overline\rho$ s'écrit
\[
\frac{b^2}{c}=\frac{\overline b^2}{\overline c}\quad,
\]
qui n'implique pas l'isom. des deux extensions ! P. ex. en car. $\ne2$, si on prend $b=\overline b=0$ la condition est vérifiée \deleted{($\rho=-2$} (il est vrai que l'on est dans le cas exceptionnel $\rho=\unsure{-|+}2$) mais l'on n'a d'isom entre les \deleted{\unsure{rev.|ext.}} $\unsure{\xi|\mathcal E|\varepsilon}_H$ et $\xi_{\overline H}$ que si $\overline c/c$ est un carré. \deleted{\unsure{et\ldots}} Etc\ldots

Il serait intéressant de voir si tout $\rho_0$ inv. peut être réalisé loc. \uuline{Zar} (et non seulement fpqc, ou fppf) Je soupçonne que oui, en partant d'un $H$, et en prenant $H'=\operatorname{int}(u)H$ pour $u\in A^{\unsure{\times|*}}$ convenable\ldots.

\uline{NB} Vérification immédiate dans le cas où $A$ est scindée\ldots
% archive sheet number: 115
\scanpage{117}
Je vais donner une expression de $\rho$. D'abord en termes d'une base
\[
1,u,v\qquad\text{base de } B=H+H'
\]
posant
\[
\begin{aligned}
&\left\{\begin{aligned} b&=\operatorname{Tr}u,\ c=\det u\\ b'&=\operatorname{Tr}v,\ c'=\det v\\ t&=\operatorname{Tr}uv=\operatorname{Tr}vu
\end{aligned}\right. &&
\begin{aligned}\delta_H&=b^2-4c\\ \delta_{H'}&=b'^2-4c'\\ &\end{aligned}
\end{aligned}
\]
on trouve l'expression \margin{\uline{NB} la condition $H$ étale $\Longleftrightarrow\delta_H$ inv $H'$ étale $\Longleftrightarrow\delta_{H'}$ inv.}
\[
\boxed{\rho_0=\frac{P\deleted{\unsure{\delta|6}}}{Q}}
\]
où
\[
\begin{aligned}
P&=P(b,c,b',c',t)=(b^2-4c)(b'^2-4c')=\delta_H\delta_{H'}^{\unsure{*|\cdot}}\\ Q&=Q(b,c,b',c',t)=t^2-bb't+(b^2c'+b'^2c-4cc')\deleted{\unsure{\delta_{\!B}}}
\end{aligned}
\]
et d'autre part on a
\[
\begin{aligned}
Q&=\unsure{-|{\deleted{-}}}\tfrac12\delta_{\varphi|B}\,\substack{\text{pour la base}\\1,u,v} =\unsure{-|{\deleted{-}}}\tfrac12\det\begin{pmatrix} \operatorname{Tr}1 & \operatorname{Tr}u & \operatorname{Tr}v\\ \operatorname{Tr}u & \operatorname{Tr}u^2 & \operatorname{Tr}uv\\ \operatorname{Tr}v & \operatorname{Tr}vu & \operatorname{Tr}v^2
\end{pmatrix}\\
&=\unsure{-|{\deleted{-}}}\,\text{discr. divisé de la forme}\\ &\qquad\det_{|B}\ \text{p.r. à la base }1,u,v
\end{aligned}
\]
(matrice \inserted{de $\varphi_{|B}$} égale :
\[
\begin{pmatrix}2&b&b'\\ b&\deleted{\unsure{2}}b^2-2c&t\\ b'&t&b'^2-2c'\end{pmatrix}\qquad)
\]
donc par hypo ($H,H'$ en position générale) \deleted{donc} on a bien $Q$ inversible, donc $\rho_0$ est défini, et non nul.

Je voudrais donner à cette formule pour $\rho_0$ une forme intrinsèque, \unsure{indépendant|indépendante} du choix des bases.
% archive sheet number: 116
\scanpage{118}
Regardons d'abord un sous-fibré vect. de rg 2 quelconque, $H$, et regardons \deleted{\unsure{la}} la forme \deleted{det} quadratique
\[
\det{}_{|H}
\]
On a vu que son discriminant \inserted{$\delta_H$} est aussi égal : celui de la forme bilinéaire \inserted{sym.} $\varphi|H$ $\operatorname{Tr}uv$ (car $H\ni1$). C'est une section de
\[
\det H^{\otimes(\unsure{-|{\deleted{+}}-}2)}\simeq\deleted{\det}(H/k)^{\otimes-2}
\]

Par rapport à une base $1,u$, d'où la base $1\wedge u$ de $\det H\simeq H/k$, et $(1\wedge u)^{\otimes-2}$ de $(\det H)^{\otimes-2}$, on trouve le discriminant
\[
\det\begin{pmatrix}\operatorname{Tr}1&\operatorname{Tr}u\\\operatorname{Tr}u&\operatorname{Tr}u^2\end{pmatrix}=b^2-4c
\]
donc
\[
\delta_H=(b^2-4c)(1\wedge u)^{\otimes-2}
\]

Si on a $H,H'$ tels que $H\cap H'=k$ (sous-vect. de rang 2 en position transverse p.r. à $k.1$), et posant
\[
B=H+H'
\]
comme dessus, on a
\[
B/k\simeq H/k+H'/k
\]
d'où isom (défini au signe près, si on ne précise les conventions)
\[
\det B\simeq\underset{k}{\underline{\det k}}\otimes\det B/k\simeq \deleted{\unsure{\det}}\unsure{H/k\otimes H/k|H'/k\otimes H/k|H/k\otimes H'/k},
\]
Donc \deleted{on peut interpréter}
\[
(\det B)^{\otimes-2}\simeq(H/k)^{\otimes-2}\otimes(H'/k)^{\otimes-2}
\]
% archive sheet number: 117
\scanpage{119}
isom qui ne dépend plus de signes. \deleted{Ceci} On peut donc considérer
\[
\left\{\begin{aligned} \delta_H\delta_{H'}&\in(\det B)^{\otimes-2}\simeq D^{\otimes-2}\\ \text{de }\overline{\mathrm m}\qquad\delta_B&\in(\det B)^{\otimes-2}\simeq D^{\otimes-2}
\end{aligned}\right.
\]

Je rappelle que la \deleted{fait} condition que \deleted{\unsure{B}}\deleted{\unsure{B}}$H,H'$ soient ``en position générale'' i.e. engendrent l'\uline{algèbre} $A$, est que
\[
\begin{aligned}
\delta_B\quad&\text{est section \uline{inv} de }(\det\unsure{D|B})^{\otimes-2}\\ &\hspace{7em}\simeq D^{\otimes-2}
\end{aligned}
\]
i.e $D$ anisotrope. Ceci dit, on trouve \margin{si $H,H'$ étales sur $k$ i.e. $\delta_H\delta_{H'}$ section inversible}
\[
\boxed{\rho_0=\frac{\delta_H\delta_{H'}}{\delta_B}}
\]

Cette formule n'est définie pour le moment que si $H,H'$ sont étales /$k$ (et en position générale). On peut la prendre comme \uline{définition} de $\rho_0$ \deleted{\unsure{dès}}\inserted{dès} que $\delta_B$ inv. i.e. $H,H'$ en position générale (pas néc\textsuperscript{t} étale). On peut même définir $\rho_0$ encore plus généralement, du moment que $\delta_H\delta_{H'}$ et $\delta_B$ ne
% archive sheet number: 118
\scanpage{120}
s'annulent pas simultanément, i.e. ils engendrent $D^{\otimes-2}$ : on trouve alors pour $\rho_0$ une section de $\mathbb P^1_k$ prenant la valeur $\infty$ aux pts où $\delta_B$ s'annule, \deleted{la valeur} i.e. $H,H'$ contenus dans un $\overline{\mathrm m}$ Borel \inserted{(mais $H,H'$ étales)}, et la valeur $0$ en les pts où $H$ ou $H'$ n'est pas étale. \unsure{L}

\uline{NB} Le seul cas sur un corps où $\rho_0$ n'est pas défini, est celui où $B=B_D$ est une sous-alg. de Borel, pour droite $D\subsetA^{\natural}$ \unsure{tot.} isotr., et $H$ (disons, à moins d'échanger les rôles) est $k\oplus D$, et enfin $H'$ est \deleted{de la forme} un \deleted{vectoriel}\inserted{\unsure{2-}plan de $B$} qui contient $k$ et est distinct de $H$ sans plus.

\deleted{Cette \unsure{situat|situat.}} Deux telles situations sont conjuguées loc -- i.e. ce type de situation est standard [correspond au choix d'un Borel et d'un tore maximal dedans, i.e. \uline{couple de Killing}.]
% archive sheet number: 119
\scanpage{121}
Revenons à l'expression explicite
\[
-\delta_{B=H+H'}\;\deleted{\unsure{\delta}}=Q\deleted{\unsure{(}}=t^2-bb't+(b^2c'+b'^2c-4cc')
\]
Je dis que pour un syst des él\textsuperscript{ts}
\[
(b,c,b',c',t)\in k^5\quad,
\]
tels que cette quantité soit inversible\textsuperscript{*},
% slanted
\margin{* Sans doute inutile de supposer $\delta$ inversible, cf plus bas} il existe toujours, \inserted{\deleted{loc. fppf sur $\operatorname{Spec} k$,}} une façon de les
% slanted
\margin{on trouvera en fait une façon \uuline{canonique}} réaliser dans une algèbre d'Azumaya, par deux él\textsuperscript{ts} $u,v$, \deleted{Pour les} tels que $H=k\oplus ku$ et $H'\deleted{\unsure{\oplus}}=k\opluskv$ sont des sous-fibrés vect. de rang 2 en position relative générale. En fait, on connaît déjà $(H,u)$ \deleted{\unsure{\textit{k}}} comme $k$-alg. munie d'un gén, à iso. un. près, par $b,c$, et itou pour $(H',v')$ par $b',c'$. Si $\overbrace{(b^2-4c)(b'^2-4c')}^{\delta\delta'}$ est
% slanted
\margin{À partir d'ici, \unsure{supposons|supposé} \unsure{tacitement|$\delta$ inversible} que $k$ est un corps alg clos. Ça simplifie le langage} inversible, on trouve une quantité $\rho_0\ne0$, correspondant à \deleted{deux} une position relative dans $M_{2\,k}$ disons (à condition de faire extension de la base vers les solutions de
\[
\lambda^2-(\rho_0+2)\lambda+1=0
\]
% archive sheet number: 120
\scanpage{122}
\deleted{\unsure{Voir|Voici|Quid} P} avec deux sous-alg. \deleted{\unsure{H}}étales (\deleted{diago} scindées) $\overline H,\overline H'$ correspondant à l'inv. $\rho_0$. Quitte à faire encore ext. biquadratique, \deleted{$H\unsure{\simeq|=}\overline H,H'$} $\overline H\simeq H$, $\overline H'\simeq H'$, ou encore \deleted{\unsure{Donc}} \deleted{\unsure{Cela est}}\inserted{OPS $H\unsure{,|\ne}H'\subset A$.} Soit alors
\[
\tau=\operatorname{Tr}uv,
\]
si on avait $t=\tau$, on aurait gagné. On sait que
\[
\deleted{b^2-4\unsure{c|\cdot}}\frac{\delta\delta'}{Q(t)}=\frac{\delta\delta'}{Q(\tau)}\qquad(=\rho_0)
\]
donc
\[
Q(t)=Q(\tau)\deleted{\unsure{=}}\qquad(=\sigma)
\]
\deleted{\unsure{Or les|On les}} ou encore, posons \deleted{Supposons $k$ corps alg. clos. Comme}
\[
\begin{aligned}
\deleted{\unsure{t^2}}\quad R(T)&=T^2-bb'T+r\\ \text{où}\qquad r&=(b^2c'+b'^2c-4cc'-\sigma)
\end{aligned}
\]
$t$ et $\tau$ sont des solutions de l'équation $R(T)=0$. Si donc on avait $\tau\ne t$, on aurait
\[
t+\tau=\deleted{\unsure{+}}bb'\qquad\text{donc}\quad\tau=bb'-t
\]
Or considérons
\[
u'=\operatorname{Tr}u\cdot1-u=b-u
\]
% archive sheet number: 121
\scanpage{123}
On aura
\[
\begin{aligned}
\operatorname{Tr}u'&=\operatorname{Tr}u=b\\ \deleted{\unsure{\operatorname{Tr}}}\det u'&=\det u=c\\ \operatorname{Tr}u'v&=\operatorname{Tr}(b-u)v=b\underbrace{\operatorname{Tr}v}_{b'}-\underbrace{\operatorname{Tr}uv}_{\tau}=bb'-\tau=t
\end{aligned}
\]
et on gagne, en remplaçant $u$ par $u'$.

% slanted
\margin{Il manque le cas où $\delta\delta'=0$ cf plus bas pour ce cas-là} On a donc gagné dans le cas où on est sur un corps. Or on va maintenant voir que l'on a une réalisation \uline{canonique} de \deleted{$H,H'$} $(b,c,b',c',t)$ dans une alg. d'Azumaya, Pour ceci, on regarde l'alg.
% slanted
\margin{Attention au cercle vicieux, si on veut éviter la vérification fastidieuse de l'associativité de la loi de composition de $M(b,c,b',c',t)$}
\[
M(b,c,b',c',t)
\]
définie par ailleurs, à l'aide d'une base formelle
\[
\underline{1}\quad\underline{u},\ \underline{v},\ \underline{u}\,\underline{v}
\]
et une multiplication explicite. Tout revient à prouver que celle-ci est Azumaya. Il suffit de le voir sur les fibres géom -- on est ramené au cas d'un corps de base alg. clos. Mais là, l'argument
% archive sheet number: 122
\scanpage{124}
de réalisabilité nous donne en $\overline{\mathrm m}$ temps que $A\simeq M(b,\unsure{c|b},b',c',t)$, laquelle est donc bien Azumaya. On gagne !

Mais j'ai oublié de faire la vérification (sur le \deleted{alg} corps alg clos) quand $\delta\delta'=0$, donc $H$ ou $H'$ est ``unipotente'' \deleted{\unsure{i.e. on}} (disons $H'$) i.e. \deleted{Considérons le discriminant de} \deleted{la forme} (comme on est sur $k$ alg clos) \deleted{$Q(t)$ considéré comme un polynôme en} $H'=\deleted{\unsure{k[v]/v^2}}k[V]/V^2$, $V$ une indét. \deleted{$L\subset A^{\unsure{\mathrm{s}|\mathrm{tn}}}$ tot isotr. i.e. $\det|L=0$ $\operatorname{Tr}|L=0$}

La formulation qui suit montre que pour $H,H'$ donnés, la \inserted{réponse au} problème de réalisabilité ne dépend pas du choix des \deleted{bases} éléments $u\in H$, $v\in H'$ utilisés pour faire des bases. \deleted{On peut donc} \inserted{Il y a donc} trois ``cas types''

1) $\delta$ \uline{et} $\delta'$ \deleted{\unsure{$\ne1$}}$\ne0$ \quad OPS $u^2=u,v^2=v$ (idempotents : donc $b=b'=1$, $c=c'=0$
\[
Q(t)=t^2-t=(t-1)t
\]

2) $\delta\deleted{\unsure{\text{ et }\delta'}}\ne0,\ \delta'\deleted{\unsure{\ne0}}=0$ (ou l'inverse) OPS \quad $u^2=u\quad v^2=\unsure{0|v}$ i.e. $\begin{aligned}b&=1&c&=0\\b'&=0&c'&=0\end{aligned}$

3) $\delta=\delta'=0$ \quad OPS \quad $u^2=0,v^2=0$ i.e $b=b'=c=c'=0$

Dans les cas $2^\circ)$ et $3^\circ)$ et dès qu'on prend $v^2=0$, (sans condition sur $u$) on trouve
\[
Q(t)=t^2
\]
% archive sheet number: 123
\scanpage{125}
Notre assertion de réalisabilité \inserted{``géom'' (sur $k$ alg. clos)} revient aux trois assertions suivantes

\uline{1}\textsuperscript{o}) Si $\deleted{\unsure{\rho_0}}t\in k$ \deleted{avec $\unsure{f_0|t}\notin\{0,1\}$,} alors $\exists$ projecteurs $u,v$ \inserted{(de rang 1)} \deleted{en position générale,} tels que
\[
\operatorname{Tr}uv=t
\]
[et si on suppose $t\notin\{0,1\}$, alors $u,v$ seront en ``position générale'']

2\textsuperscript{o}) Si $t\in k$, alors $\exists$ projecteur \inserted{(rg 1)} $u$, et $v$ nilpotent $\ne0$, \deleted{dont $\exists$} tels que
\[
\operatorname{Tr}uv=t
\]
(et si $t\ne0$, on aura $u,v$ en pos\unsure{.} gén.)

3\textsuperscript{o}) Si $t\in k$, alors $\exists$ $u,v$ nilp.\textsuperscript{ts} $\ne0$, tels que
\[
\operatorname{Tr}uv=t
\]
(et si $t\ne0$, $u,v$ sont en position générale)

Interprétant $A=\operatorname{End}(V)$ $V$ vect. de dim 2, on peut écrire dans ts les cas
\[
u=\deleted{\unsure{x\otimes x'}}x'\otimes x,\qquad v=\deleted{\unsure{y\otimes y'}}y'\otimes y\qquad x,y\in V,\ x',y'\in V^\vee
\]
donc
\[
\begin{aligned}
u\circ v&=\langle\deleted{\unsure{x',y'}}y,x'\rangle\, y'\otimes x\unsure{=|\cdot}\\ \unsure{\text{o}|\text{ou}}\quad t&=\operatorname{Tr}uv=\langle x,y'\rangle\langle y,x'\rangle
\end{aligned}
\]
\deleted{et il faut voir}

\uline{Cas 1} \deleted{$\unsure{\langle x,x'\rangle}=1,\langle y,y'\rangle=1$} $\langle x,x'\rangle=1,\langle y,y'\rangle=1$ On choisit base $e_1,e_2$ de $V$ tel que $x=e_1$, $x'=e'_1$. Si les coeff. de $y,y'$ sont $(y_1,y_2)$ $(y'_1,y'_2)$, la relation $\langle y,y'\rangle=1$
% archive sheet number: 124
\scanpage{126}
s'écrit
\[
y_1y'_1+y_2y'_2=1
\]
et on a
\[
t=y_1^{\deleted{\unsure{'|1}}}y'_1
\]
\deleted{On peut prendre} les relations deviennent \deleted{$y_1=1,y'_1=t$ et $y_2,y'_2$ tel que $y_1y'_1=1-t$},
\[
\left\{\begin{aligned}y_1y'_1&=t\\ y_2y'_2&=1-t\end{aligned}\right.
\]

Si $t\ne0,1$ on peut choisir arbitrairement $y_1,y_2$ \deleted{$\unsure{x|\in}$} $\in k^{\times}$, et tirer \deleted{$y'_1,\unsure{y'|y'_2}$}
\[
y'_1=t/y_1,\quad y'_2=\frac{1-t}{y'_2}
\]
Si $t=0$, \deleted{on a $y_1$ \unsure{.}} ça devient
\[
\begin{aligned} y_1y'_1&=0\\ y_2y'_2&=1\end{aligned}
\]
deux familles de solutions \inserted{à deux par.} suivant que $y_1=0$ ou $y'_1=0$ et symétriques si $t=1$.

\uline{Cas 2} On peut tjrs $x=e_1$, $x'=e'_1$, \unsure{relation|relations}
\[
\begin{aligned} y_1y'_1+y_2y'_2&=0\\ t&=y_1y'_1\end{aligned}
\]
i.e. comme équations en les $y_i,y'_i$ : \deleted{$\unsure{y_2}$}
\[
\left\{\begin{aligned} y_1y'_1&=t\\ y_2\unsure{y'_2|y_2^2}&=-t\end{aligned}\right.
\]

\uline{Cas 3} On peut prendre $x=e_1$, $x'=e'_{\underline{\underline 2}}$.
\[
\deleted{\unsure{y_1y'_1+y_2y'_2=0}}
\]
On doit avoir
\[
\left\{\begin{aligned} y_1y'_1+y_2y'_2&=0\\ t&=y'_1y_2\end{aligned}\right.
\]
\deleted{\unsure{et}} \ldots
% archive sheet number: 125
\scanpage{127}
\uline{NB} Finalement, on n'a pas eu besoin du tout de faire intervenir $\rho$ dans la vérification, c'était une complication inutile

\begin{center}\rule{3cm}{0.4pt}\end{center}

Je vais reprendre la théorie précédente. Soient $H_{\deleted{\unsure{k}}},H'$ deux fibrés vectoriels de rang 2 \unsure{sur|\deleted{sur}} $k$, munis d'une \deleted{forme $q$} structure de $k$-algèbre (néc. comm.) d'où \deleted{forme}\inserted{sur} $H$ et $H'$ formes \uline{norme} et trace
\diagram{\begin{tikzcd}[ampersand replacement=\&,column sep=small,row sep=small]
{\deleted{\Delta}N_H:H}\arrow[r]\& k\&{\deleted{\Delta}N_{H'}:H'}\arrow[r,dash]\&k\\ {T_H:H}\arrow[r]\&k\&{T_{H'}:H'}\arrow[r,dash]\&k\\ {\varepsilon_H:\deleted{\unsure{1_H}}k}\arrow[r]\&H\&{\varepsilon_{H'}:k'}\arrow[r,dash]\&H'
\end{tikzcd}}
% slanted
\margin{$\varepsilon_H(\lambda)=\lambda1_H$ \quad $\varepsilon_{H'}(\lambda)=\unsure{\rho_0|\deleted{\rho_0}}\lambda.1_{H'}$} \deleted{l'essentiel de ce qui va suivre, n'aura pas besoin de la structure multiplicative, mais seulement des formes précédentes (voire, seulement de $N$ et $T_H$)} La formule
% slanted
\margin{\uline{NB} La forme $T_H$ se déduit de la donnée de $1_H\in H$ (i.e. $\varepsilon_H$) et de $\Delta_H$, comme $T_H(u)=\varphi_H(u,1)$ et le fait que $\Delta_H(1)=1$ (d'où $T(1)=2$) implique que les f\textsuperscript{es} $\varphi_H(1,u)=\psi_H(1,u,)\quad\forall u$}
\[
\underbrace{\varphi_{\deleted{\unsure{N|\Delta}}H}}_{=\varphi_{\Delta_H}}(u,v)=\operatorname{Tr}(u)\operatorname{Tr}(v)-\underbrace{\operatorname{Tr}(uv)}_{\psi_H(u,v)}
\]
montre que la forme $\operatorname{Tr}(uv)$ sur $H$ s'exprime par la connaissance de $N_H$ et $T_H$ comme
\[
\operatorname{Tr}_{H}(uv)=T_H(u)T_H(v)-\varphi_{\deleted{\Delta}H}(u,v)
\]
i.e.
\[
\varphi_H(u,v)+\psi_H(u,v)=T_H(u)T_H(v)
\]
% archive sheet number: 126
\scanpage{128}
\unsure{Posons|poser}
\[
B=H\mathbin{\amalg}_k H'=H\oplus H'/\deleted{\unsure{k1_H+k1_{H'}|k(1_H+1_{H'})}}\inserted{k(1_H\ominus1_{H'})}
\]
% slanted
\margin{Mais la structure d'alg. de $H$ résulte déjà de sa structure de module quadratique, ``muni de $1_H$ tel que $\Delta_H(1)=1$'', et itou pour $H'$ !} Ainsi $H$ et $H'$ se réalisent comme sous-fibrés vect. de rg 2 d'un $B$, et
\[
H\cap H'=k1_B\qquad H+H'=B
\]
où $1_B=1_H=1_{H'}$ est l'image commune de $1_H$ et $1_{H'}$ dans $B$.

On considère aussi
\[
P=H\otimes_k H'
\]
et les plong\textsuperscript{ts}
\diagram{\begin{tikzcd}[ampersand replacement=\&,column sep=small,row sep=small]
\&P\&\&[-1em] \\ H\arrow[ur,"i_0"]\& \& H'\arrow[ul,"j_0"']\& {\begin{gathered}i(u)=u\otimes1\\ j(v)=1\otimes v\end{gathered}}
\end{tikzcd}}
On peut alors réaliser $B$ comme sous-fibré vectoriel de $P$, \deleted{\illegible{1}}\inserted{car} $i,j$ [\deleted{\unsure{le|lequel}} qui ``coïncident sur $k$'' i.e. $i(\lambda1_H)=j(\lambda1_{H'})=\lambda\underbrace{1_H\otimes1_{H'}}_{\text{``est'' }1_P}\bigr]$ se factorisent de façon unique par $B\longrightarrow P$
\diagram{\begin{tikzcd}[ampersand replacement=\&,column sep=small,row sep=small]
\&P\&\\ \&B\arrow[u]\arrow[dl]\arrow[dr]\&\\ H\arrow[uur,"i_0"]\&\&H'\arrow[uul,"j_0"']
\end{tikzcd}}
% slanted
\margin{Attention, \unsure{Priori pas|$B$ n'a pas|$P$ n'a pas} de structure d'algèbre !} \margin{NB on a $B\xrightarrow{T_B}k$ induisant $T_H,T_{H'}$}
% slanted
\margin{(puisque $T_H(1_H)=T_{H'}(1_{H'})=2$)}
% Faint mirror-image typescript shows through the sheet.
% archive sheet number: 127
